\documentclass[11pt]{article}

\usepackage[margin=1in]{geometry}
\usepackage[T1]{fontenc}
\usepackage[utf8]{inputenc}
\usepackage{lmodern}
\usepackage{amsmath,amssymb,amsthm,mathtools}
\usepackage{enumitem}
\usepackage{booktabs}
\usepackage{xcolor}
\usepackage{hyperref}
\usepackage[nameinlink,capitalize]{cleveref}
\usepackage{microtype}

\hypersetup{colorlinks=true,linkcolor=blue,citecolor=blue,urlcolor=blue}
\emergencystretch=3em

\newtheorem{theorem}{Theorem}[section]
\newtheorem{lemma}[theorem]{Lemma}
\newtheorem{proposition}[theorem]{Proposition}
\newtheorem{corollary}[theorem]{Corollary}
\newtheorem{conjecture}[theorem]{Conjecture}
\newtheorem{problem}[theorem]{Problem}
\theoremstyle{definition}
\newtheorem{definition}[theorem]{Definition}
\newtheorem{remark}[theorem]{Remark}
\newtheorem{audit}[theorem]{Audit note}

\newcommand{\E}{\mathbb E}
\newcommand{\Pp}{\mathbb P}
\newcommand{\R}{\mathbb R}
\newcommand{\N}{\mathbb N}
\newcommand{\one}{\mathbf 1}
\newcommand{\Pois}{\operatorname{Pois}}
\newcommand{\Var}{\operatorname{Var}}
\newcommand{\Cov}{\operatorname{Cov}}
\newcommand{\supp}{\operatorname{supp}}
\newcommand{\DeltaN}{\Delta_{n-1}}
\newcommand{\GT}{\operatorname{GT}}
\newcommand{\diag}{\operatorname{diag}}

\title{Conjecture R2 for Poissonized Generalized Good--Turing Risk\\
\large Active-uniform and Dirac--uniform least-favorable laws}
\author{Working research draft}
\date{\today}

\begin{document}
\maketitle

\begin{abstract}
We audit and prove a precise version of Conjecture R2.  Under Poissonization, for the canonical order-$r$ Good--Turing estimator of the order-$r$ occupancy mass, the squared risk is a separable simplex objective.  After the scaling $z_i=t p_i$, its coordinate kernel is
\[
  \varphi_\theta(z)=z^\theta(z+\theta)e^{-z},
  \qquad \theta=r+1.
\]
We prove that every finite-alphabet least-favorable distribution is uniform on its active support or Dirac--uniform, i.e. it has at most one exceptional positive atom and all other positive atoms are equal.  The proof is a crossing-plus-structured-perturbation argument: the marginal equation $\varphi_\theta'(z)=\lambda$ has at most two positive roots, and any root in a convex branch can occur with multiplicity at most one.  This gives the desired active-uniform structure for every order $r\ge0$.  The audit identifies two necessary wording corrections.  First, for $r\ge1$ the exceptional atom can be light as well as heavy, so the invariant statement is one-exceptional or Dirac--uniform, not literally one-heavy.  Second, for $r\ge1$ boundary zeros are allowed by the KKT inequalities and should be interpreted as reducing the active support.  We also record the exact finite-dimensional reduction, links with the active-uniform program for missing-mass functionals, and a set of related open problems for phase diagrams, de-Poissonization, biased generalized Good--Turing estimators, and countable alphabets.
\end{abstract}

\tableofcontents

\section{Precise statement of R2}

Let $p=(p_1,\ldots,p_n)$ be a distribution on a finite alphabet, and let
\[
  N_i\sim\Pois(t p_i),\qquad i=1,\ldots,n,
\]
independently.  Here $t>0$ is the Poisson sampling intensity.  For an integer $r\ge0$, define the order-$r$ occupancy mass
\[
  M_r=\sum_{i=1}^n p_i\one\{N_i=r\}
\]
and the canonical Poissonized Good--Turing estimator
\[
  \widehat M_r=\frac{r+1}{t}\sum_{i=1}^n \one\{N_i=r+1\}.
\]
This is the Poissonized analogue of estimating the total probability mass of symbols observed exactly $r$ times.  The case $r=0$ is the usual missing-mass Good--Turing estimator.

\begin{definition}[Active-uniform and Dirac--uniform laws]
A distribution $p\in\Delta_{n-1}$ is \emph{active-uniform} if, after deleting zero coordinates, all remaining coordinates are equal.  It is \emph{Dirac--uniform} or \emph{one-exceptional} if, after a permutation,
\[
  p=(\underbrace{a,\ldots,a}_{k-1\text{ times}},b,0,\ldots,0)
\]
for some $1\le k\le n$, $a>0$, and $b\ge0$, with the convention that $k=1$ is a Dirac mass and $b=a$ is the active-uniform case.  If $b>a$ we call the exceptional atom heavy; if $0<b<a$ we call it light.
\end{definition}

The audited version of Conjecture R2 is the following theorem.

\begin{theorem}[Conjecture R2, canonical Poissonized form]\label{thm:R2}
Fix $n\ge1$, $t>0$, and $r\in\{0,1,2,\ldots\}$.  Every maximizer over $\Delta_{n-1}$ of the Poissonized squared risk
\[
  R_{r,t}(p)=\E_p\bigl(\widehat M_r-M_r\bigr)^2
\]
is active-uniform or Dirac--uniform.  Equivalently, every least-favorable finite-alphabet distribution has, after a permutation, the form
\[
  (\underbrace{a,\ldots,a}_{k-1\text{ times}},b,0,
  \ldots,0).
\]
For $r=0$, a nonuniform full-support maximizer has one heavy coordinate.  For $r\ge1$, the exceptional coordinate need not be heavy; light-exceptional local maxima are not excluded by the structural theorem.
\end{theorem}

\begin{audit}[Scope of the theorem]
The theorem proves R2 for the \emph{canonical unbiased order-$r$ Good--Turing family}.  The phrase ``generalized Good--Turing'' is also used in the literature for biased or two-frequency estimators of missing mass, for instance estimators involving both $\Phi_1$ and $\Phi_2$.  Those risks generally have a vector-separable form, with a squared bias term, rather than the purely separable kernel treated here.  The same frozen-kernel method gives finite-level criteria for such estimators, but the Dirac--uniform conclusion requires additional one-valley or root-count hypotheses.  See \cref{sec:biased-generalized}.
\end{audit}

\section{Poissonized risk reduction}

The key simplification is coordinatewise unbiasedness under Poissonization.

\begin{lemma}[Coordinatewise unbiasedness and separability]\label{lem:risk-separable}
For each coordinate $i$,
\[
  \E\left[\frac{r+1}{t}\one\{N_i=r+1\}\right]
  =\E\left[p_i\one\{N_i=r\}\right].
\]
Consequently,
\[
  R_{r,t}(p)
  =\sum_{i=1}^n
  \left[
  p_i^2 e^{-t p_i}\frac{(t p_i)^r}{r!}
  +\left(\frac{r+1}{t}\right)^2
  e^{-t p_i}\frac{(t p_i)^{r+1}}{(r+1)!}
  \right].
\]
Equivalently, with
\[
  z_i=t p_i,
  \qquad \theta=r+1,
\]
we have
\[
  R_{r,t}(p)=\frac{1}{t^2 r!}
  \sum_{i=1}^n \varphi_\theta(z_i),
  \qquad
  \varphi_\theta(z)=z^\theta(z+\theta)e^{-z},
\]
under the simplex constraint
\[
  z_i\ge0,
  \qquad \sum_{i=1}^n z_i=t.
\]
\end{lemma}

\begin{proof}
Since $N_i\sim\Pois(t p_i)$,
\[
  \E\left[\frac{r+1}{t}\one\{N_i=r+1\}\right]
  =\frac{r+1}{t}e^{-t p_i}\frac{(t p_i)^{r+1}}{(r+1)!}
  =p_i e^{-t p_i}\frac{(t p_i)^r}{r!}
  =\E\left[p_i\one\{N_i=r\}\right].
\]
The coordinate errors are independent and have mean zero.  Moreover, the events $\{N_i=r\}$ and $\{N_i=r+1\}$ are disjoint, so the coordinate squared error is
\[
  p_i^2\one\{N_i=r\}
  +\left(\frac{r+1}{t}\right)^2\one\{N_i=r+1\}.
\]
Taking expectations and summing gives the displayed expression.  The change of variables $z_i=t p_i$ gives
\[
  p_i^2\frac{(t p_i)^r}{r!}
  +\left(\frac{r+1}{t}\right)^2\frac{(t p_i)^{r+1}}{(r+1)!}
  =\frac{1}{t^2r!}\bigl(z_i^{r+2}+(r+1)z_i^{r+1}\bigr),
\]
which is the claimed kernel with $\theta=r+1$.
\end{proof}

Thus the structural problem is independent of the prefactor $1/(t^2r!)$.

\begin{definition}[Scaled risk problem]\label{def:scaled}
For $\theta\ge1$ define
\[
  \mathcal R_{\theta,t,n}(z)=\sum_{i=1}^n\varphi_\theta(z_i),
  \qquad
  \varphi_\theta(z)=z^\theta(z+\theta)e^{-z},
\]
on the compact simplex
\[
  \Sigma_{n,t}=\{z\in[0,\infty)^n:\sum_i z_i=t\}.
\]
A maximizer of $\mathcal R_{\theta,t,n}$ corresponds exactly to a least-favorable distribution for order $r=\theta-1$ by $p_i=z_i/t$.
\end{definition}

\section{The marginal curve and its root count}\label{sec:marginal}

The proof of R2 rests on a complete qualitative description of
\[
  q_\theta(z):=\varphi_\theta'(z).
\]

\begin{lemma}[First and second derivatives]\label{lem:derivatives}
For $\theta\ge1$ and $z>0$,
\[
  \varphi_\theta'(z)=z^{\theta-1}e^{-z}(\theta^2+z-z^2),
\]
and
\[
  \varphi_\theta''(z)=z^{\theta-2}e^{-z}P_\theta(z),
\]
where
\[
  P_\theta(z)=z^3-(\theta+2)z^2-\theta(\theta-1)z+\theta^2(\theta-1).
\]
The unique positive zero of $\theta^2+z-z^2$ is
\[
  \rho_\theta=\frac{1+\sqrt{1+4\theta^2}}{2}.
\]
Thus $q_\theta(z)>0$ on $(0,\rho_\theta)$ and $q_\theta(z)<0$ on $(\rho_\theta,\infty)$.
\end{lemma}

\begin{proof}
Differentiate
\[
  \varphi_\theta(z)=z^{\theta+1}e^{-z}+\theta z^\theta e^{-z}.
\]
The displayed formulae follow by collecting powers of $z$.  The quadratic $\theta^2+z-z^2$ has exactly one positive zero, namely $\rho_\theta$, and its sign gives the sign of $q_\theta$ because $z^{\theta-1}e^{-z}>0$.
\end{proof}

\begin{lemma}[Curvature sign pattern]\label{lem:curvature}
For $\theta=1$,
\[
  \varphi_1''(z)=z(z-3)e^{-z};
\]
hence $q_1$ is strictly decreasing on $(0,3)$ and strictly increasing on $(3,\infty)$.

For every $\theta>1$, the polynomial $P_\theta$ has exactly two positive zeros
\[
  0<\alpha_\theta<\rho_\theta<\beta_\theta,
\]
and
\[
  \varphi_\theta''>0\text{ on }(0,\alpha_\theta),
  \qquad
  \varphi_\theta''<0\text{ on }(\alpha_\theta,\beta_\theta),
  \qquad
  \varphi_\theta''>0\text{ on }(\beta_\theta,\infty).
\]
Thus, for $\theta>1$, $q_\theta$ increases, then decreases, then increases.
\end{lemma}

\begin{proof}
The case $\theta=1$ follows immediately from \cref{lem:derivatives}, since then
\[
  P_1(z)=z^2(z-3)
\]
and $z^{\theta-2}P_\theta(z)=z^{-1}z^2(z-3)=z(z-3)$.

Assume now $\theta>1$.  The signs of the coefficients of $P_\theta$ are
\[
  +,\quad -,\quad -,\quad +,
\]
so Descartes' rule of signs gives at most two positive zeros.  On the other hand,
\[
  P_\theta(0)=\theta^2(\theta-1)>0,
  \qquad
  \lim_{z\to\infty}P_\theta(z)=+\infty.
\]
At $z=\rho_\theta$, use $\rho_\theta^2=\rho_\theta+\theta^2$ to obtain
\[
  P_\theta(\rho_\theta)=-\rho_\theta-2\theta^2<0.
\]
Hence there is at least one positive zero below $\rho_\theta$ and at least one above it.  Since there are at most two, there are exactly two, denoted $\alpha_\theta$ and $\beta_\theta$.  The sign pattern follows from $P_\theta(0)>0$, $P_\theta(\rho_\theta)<0$, and $P_\theta(z)\to+\infty$.
\end{proof}

\begin{lemma}[Two-root marginal equation]\label{lem:two-root}
For every $\theta\ge1$ and every $\lambda\in\R$, the equation
\[
  q_\theta(z)=\lambda,
  \qquad z>0,
\]
has at most two positive solutions.

More precisely, for $\theta>1$ there are two possible two-root configurations.  If $\lambda>0$, then one root lies in the convex interval $(0,\alpha_\theta)$ and the other lies in the concave interval $(\alpha_\theta,\rho_\theta)$.  If $\lambda<0$, then one root lies in the concave interval $(\rho_\theta,\beta_\theta)$ and the other lies in the convex interval $(\beta_\theta,\infty)$.  For $\lambda=0$, the only positive root is $\rho_\theta$.
\end{lemma}

\begin{proof}
For $\theta=1$, $q_1$ decreases on $(0,3)$ and increases on $(3,\infty)$ by \cref{lem:curvature}; hence a horizontal line intersects the graph at most twice.

Let $\theta>1$.  By \cref{lem:curvature}, $q_\theta$ increases on $(0,\alpha_\theta)$, decreases on $(\alpha_\theta,\beta_\theta)$, and increases on $(\beta_\theta,\infty)$.  Also
\[
  q_\theta(0+)=0,
  \qquad q_\theta(\rho_\theta)=0,
  \qquad q_\theta(\infty)=0^-.
\]
Since $\alpha_\theta<\rho_\theta<\beta_\theta$, we have $q_\theta(\alpha_\theta)>0$ and $q_\theta(\beta_\theta)<0$.  Therefore positive levels can be hit only once on $(0,\alpha_\theta)$ and once on $(\alpha_\theta,\rho_\theta)$; negative levels can be hit only once on $(\rho_\theta,\beta_\theta)$ and once on $(\beta_\theta,\infty)$; and the zero level has the unique positive root $\rho_\theta$.  This proves the claim.
\end{proof}

\section{Structured perturbations on simplex faces}

We now record the elementary local-maximum test that turns the root count into a one-exceptional conclusion.

\begin{lemma}[KKT equations on the active face]\label{lem:kkt}
Let $z\in\Sigma_{n,t}$ be a local maximum of $\mathcal R_{\theta,t,n}$, and let
\[
  A=\{i:z_i>0\}
\]
be its active support.  If $|A|\ge2$, then there is a scalar $\lambda$ such that
\[
  q_\theta(z_i)=\lambda,
  \qquad i\in A.
\]
If $j\notin A$, then the boundary KKT inequality is
\[
  q_\theta(0+)\le \lambda.
\]
In particular, when $\theta=1$ no zero coordinate can occur at a full-simplex local maximum with at least two active coordinates, while for $\theta>1$ a boundary local maximum can occur only when $\lambda\ge0$.
\end{lemma}

\begin{proof}
On the relative interior of the face $A$, the only equality constraint is $\sum_{i\in A}z_i=t$.  The Lagrange multiplier rule gives $q_\theta(z_i)=\lambda$ for all active coordinates.  For an inactive coordinate $j$ and active coordinate $i$, move mass $\varepsilon$ from $i$ to $j$.  The first-order change is
\[
  \varphi_\theta(\varepsilon)-\varphi_\theta(0)
  +\varphi_\theta(z_i-\varepsilon)-\varphi_\theta(z_i)
  =\varepsilon\bigl(q_\theta(0+)-q_\theta(z_i)\bigr)+o(\varepsilon),
\]
when the one-sided derivative $q_\theta(0+)$ is finite.  For $\theta=1$, $q_1(0+)=1$, and $q_1(z)<1$ for all $z>0$, so the inequality fails unless there is only one active coordinate.  For $\theta>1$, $q_\theta(0+)=0$, so the inequality is $0\le\lambda$.
\end{proof}

\begin{lemma}[Convex-branch singleton test]\label{lem:singleton}
Let $z\in\Sigma_{n,t}$ be a local maximum of $\mathcal R_{\theta,t,n}$ on its active face.  If two active coordinates are equal to $x>0$, then
\[
  \varphi_\theta''(x)\le0.
\]
Consequently, an active value lying in a strictly convex branch of $\varphi_\theta$ can occur with multiplicity at most one.
\end{lemma}

\begin{proof}
If $z_i=z_j=x$, use the feasible balanced perturbation
\[
  (x,x)\mapsto(x+\varepsilon,x-\varepsilon).
\]
For small $\varepsilon$, this stays in the same active face and preserves the sum.  The second-order change is
\[
  \varphi_\theta(x+\varepsilon)+\varphi_\theta(x-\varepsilon)-2\varphi_\theta(x)
  =\varphi_\theta''(x)\varepsilon^2+o(\varepsilon^2).
\]
If $\varphi_\theta''(x)>0$, the objective increases for sufficiently small nonzero $\varepsilon$, contradicting local maximality.
\end{proof}

\section{Proof of R2}

\begin{proof}[Proof of \cref{thm:R2}]
By \cref{lem:risk-separable}, it is enough to maximize $\mathcal R_{\theta,t,n}$ on $\Sigma_{n,t}$, with $\theta=r+1\ge1$.  A maximizer exists by compactness.  If it has active support size one, it is a Dirac law and hence Dirac--uniform.  Assume from now on that the active support has size at least two.

By \cref{lem:kkt}, all positive active coordinates solve the same scalar equation
\[
  q_\theta(z)=\lambda.
\]
By \cref{lem:two-root}, there are at most two positive active values.  If there is only one value, the maximizer is uniform on its active support.

Suppose there are two positive active values.  By \cref{lem:two-root}, exactly one of them lies in a convex branch of $\varphi_\theta$: for $\lambda>0$ it is the smaller value in $(0,\alpha_\theta)$, while for $\lambda<0$ it is the larger value in $(\beta_\theta,\infty)$.  For $\theta=1$, the only two-root case is the negative-level case, and the convex root is the larger value $>3$.  By \cref{lem:singleton}, the convex-branch value can occur with multiplicity at most one.  Hence all remaining active coordinates are equal, and the maximizer is one-exceptional on its active support.  Reintroducing zero coordinates gives the asserted Dirac--uniform form for $p_i=z_i/t$.

For $r=0$, i.e. $\theta=1$, the two-root case has one repeated value in $(\rho_1,3)$ and one singleton value in $(3,\infty)$, so the exceptional coordinate is heavy.  For $r\ge1$, i.e. $\theta>1$, the positive-level two-root case gives a singleton light value and a repeated larger value below $\rho_\theta$, while the negative-level case gives a singleton heavy value.  This proves the final wording claim.
\end{proof}

\begin{audit}[What changed after audit]
The first-pass proof idea is sound for the separable canonical order-$r$ family, but it requires the explicit marginal-root analysis in \cref{lem:two-root}.  The following corrections are important.
\begin{enumerate}[label=\textup{(\roman*)}]
\item Descartes' rule applies to the curvature polynomial $P_\theta$, not directly to the KKT equation $q_\theta(z)=\lambda$.  One must combine the curvature sign pattern with $q_\theta(0+)$, $q_\theta(\rho_\theta)=0$, and $q_\theta(\infty)=0^-$.
\item For $r\ge1$, the statement should be ``one-exceptional'' or ``Dirac--uniform,'' not literally ``one-heavy.''  A positive KKT level has a singleton light value.
\item For $r\ge1$, $q_\theta(0+)=0$, so boundary zeros are compatible with local maximality whenever the active multiplier satisfies $\lambda\ge0$.  Thus the clean statement is active-support uniformity/Dirac--uniformity, with zeros allowed outside the active support.
\end{enumerate}
\end{audit}

\section{Finite-dimensional reduction of the least-favorable risk}\label{sec:finite-reduction}

The structural theorem reduces the finite-alphabet least-favorable value to a small family of one-dimensional searches.

\begin{proposition}[Exact stratum reduction]\label{prop:stratum}
For $\theta\ge1$, $t>0$, and $n\ge1$,
\[
  \max_{z\in\Sigma_{n,t}}\sum_i\varphi_\theta(z_i)
  =\max_{1\le k\le n}\;\max_{0\le b\le t}
  \left\{(k-1)\varphi_\theta\!
    \left(\frac{t-b}{k-1}\right)+\varphi_\theta(b)\right\},
\]
where for $k=1$ the inner expression is interpreted as $\varphi_\theta(t)$.  The active-uniform candidate of support size $k$ is included by taking $b=t/k$.
\end{proposition}

\begin{proof}
By \cref{thm:R2}, some maximizer has Dirac--uniform form with active support size $k$ and exceptional scaled mass $b$.  If $k\ge2$, the remaining $k-1$ positive coordinates are all equal to $(t-b)/(k-1)$.  Conversely, every such vector belongs to $\Sigma_{n,t}$ after adding $n-k$ zeros, so comparing these strata is exact.
\end{proof}

\begin{remark}[Stationarity on a stratum]
For $k\ge2$, an interior exceptional point satisfies
\[
  q_\theta(b)=q_\theta\left(\frac{t-b}{k-1}\right).
\]
This scalar equation is useful for numerical phase diagrams.  The theorem above does not assert that the global maximizer is always full-support, nor that the exceptional atom is always heavy; both questions depend on $r$, $t$, and $n$.
\end{remark}

\section{Relation to the active-uniform program}\label{sec:connections}

The proof is an instance of the active-uniform method developed for missing-mass functionals.  The general template is:
\begin{quote}
\emph{KKT crossing controls possible levels, and balanced perturbations rule out repeated convex levels.}
\end{quote}
For separable objectives $\sum_i f(p_i)$, the KKT equation is one-dimensional:
\[
  f'(p_i)=\lambda.
\]
If this equation has few roots and convex roots are unstable under
\[
  (x,x)\mapsto(x+\varepsilon,x-\varepsilon),
\]
then local maxima have few active levels.  In the classical missing-mass expectation problem, the kernel $x(1-x)^t$ has a one-valley derivative, giving uniform or one-heavy extrema.  In the ordinary Poissonized Good--Turing MSE case $r=0$, the kernel after scaling is
\[
  \varphi_1(z)=z(z+1)e^{-z},
\]
which is precisely the base case of the present theorem.

The novelty in R2 is that for every $r\ge0$ the risk kernel remains just simple enough: although the curvature polynomial becomes cubic, its sign pattern still forces the marginal equation to have at most two positive roots.  This keeps the least-favorable shape in the active-uniform/Dirac--uniform class.

\subsection{Connection with high-order occupancy and high-order missing mass}

The order-$r$ occupancy mass
\[
  M_r=\sum_i p_i\one\{N_i=r\}
\]
should be distinguished from the high-order missing-mass functional sometimes written
\[
  \Theta_s=\sum_i p_i^s\one\{N_i=0\}.
\]
Both are missing- or rare-species functionals, but their risk geometries differ.  In the present order-$r$ Good--Turing setting, the Poissonized unbiased estimator produces the separable kernel $z^{r+1}(z+r+1)e^{-z}$.  For high-order missing mass, expectation kernels have the form $p^s e^{-t p}$ or $p^s(1-p)^t$, whose derivatives can have more turning points; the active-level theorem usually gives bounded-level complexity, not automatically a Dirac--uniform theorem.

\subsection{Connection with vector-separable MSE objectives}\label{sec:biased-generalized}

For a general symmetric coordinate estimator $T_i(N_i)$ of a target coordinate $A_i(p_i,N_i)$, the Poissonized MSE has the form
\[
  \operatorname{MSE}(p)
  =\sum_i v(p_i)+\left(\sum_i e(p_i)\right)^2,
\]
where $e(x)$ is the coordinate bias and $v(x)$ is the coordinate variance.  This is vector-separable, not purely separable unless the estimator is coordinatewise unbiased.  At a regular local maximum, the frozen marginal equation is
\[
  v'(x)+2B e'(x)=\lambda,
  \qquad B=\sum_i e(p_i),
\]
and a repeated level must satisfy
\[
  v''(x)+2B e''(x)\le0.
\]
Thus a Painsky-style or otherwise biased one-parameter generalized Good--Turing family will inherit R2 only if every relevant frozen kernel has the same two-root/one-convex-root geometry.  This gives a concrete route to proving or disproving the broader biased-estimator version of R2.

\section{Related open problems and conjectures}

\begin{problem}[Phase diagram for order-$r$ Poissonized GT risk]
For fixed $r\ge1$ and $n$, determine the full phase diagram in $t$ after the stratum reduction of \cref{prop:stratum}.  In particular, decide when the winner is active-uniform, one-heavy, one-light, or supported on a proper face.  The structural theorem reduces the problem to finitely many one-dimensional comparisons, but does not by itself give a single threshold.
\end{problem}

\begin{problem}[Asymptotic thresholds for $r\ge1$]
For fixed $r\ge1$ and $n\to\infty$, derive the analogue of the golden-ratio expansion known for $r=0$.  A natural surrogate is obtained by placing the repeated coordinates near the maximizer of $\varphi_{r+1}$, namely $\rho_{r+1}$, and comparing
\[
  n\varphi_{r+1}(t/n)
  \quad\text{with}\quad
  (n-1)\varphi_{r+1}(\rho_{r+1}).
\]
The audit suggests that a light-exceptional branch may also be relevant for some ranges of $t$.
\end{problem}

\begin{conjecture}[Canonical order-$r$ single-crossing, qualified]
For every fixed $r\ge0$ and $n\ge2$, each active support size $k$ has at most one value transition between its active-uniform branch and its nonuniform active branch.  The global phase diagram is then a finite ordered sequence of support and one-exceptional transitions.
\end{conjecture}

\begin{problem}[Biased generalized Good--Turing families]
Consider a one-parameter estimator using several frequencies of frequencies, for example
\[
  \widehat M_{0,a}=c_1(a)\Phi_1+c_2(a)\Phi_2.
\]
Under Poissonization, compute the functions $v_a(x)$ and $e_a(x)$ and classify the frozen kernels
\[
  \ell_{a,B}'(x)=v_a'(x)+2B e_a'(x).
\]
Find the exact parameter range for which every frozen marginal equation has at most two positive roots with at most one convex root.  In that range R2 extends; outside it one should search for three-level counterexamples.
\end{problem}

\begin{problem}[De-Poissonization]
For fixed sample size, multinomial counts are dependent and the risk is no longer a simple sum of coordinate risks.  Determine whether the Poissonized Dirac--uniform least-favorable structure survives exactly, asymptotically, or only as an approximate extremal principle.
\end{problem}

\begin{problem}[Countable alphabets and nonattainment]
Without a finite alphabet bound, least-favorable laws may fail to exist.  For $r=0$, splitting mass into many tiny atoms is a familiar obstruction.  Determine, for each $r\ge0$, whether the supremum over countable alphabets is attained, and if not, whether it is approached by active-uniform finite-support laws whose support tends to infinity.
\end{problem}

\begin{conjecture}[Frozen-kernel active-support principle]
In constrained, biased, or weakly nonseparable variants of missing-mass risk, every regular least-favorable law belongs to a finite active-level stratum determined by the roots of a frozen one-dimensional marginal equation.  Uniform and Dirac--uniform laws are the universal two-root cases; higher-level laws arise only when the frozen marginal has higher crossing number.
\end{conjecture}

\section{Bibliographic notes}

Good--Turing estimators originate in Good's account of Turing's wartime frequency-estimation method \cite{Good1953}.  McAllester and Schapire gave influential learning-theoretic convergence guarantees for Good--Turing missing-mass estimation \cite{McAllesterSchapire2000}.  Berend and Kontorovich developed sharp missing-mass expectation bounds and finite-alphabet extremizer structure \cite{BerendKontorovich2012}.  Painsky's JASA paper introduced a generalized Good--Turing framework for missing-mass estimation and optimized MSE upper bounds over its parameters \cite{Painsky2023}; his JMLR paper studies convergence guarantees for Good--Turing occupancy-probability estimators and discusses generalized Good--Turing variants using additional frequencies of frequencies \cite{Painsky2022}.  Favaro and Naulet study high-order missing masses and the rare-type match problem \cite{FavaroNaulet2023}.  These works motivate the distinction between the canonical unbiased order-$r$ Poissonized theorem proved here and the broader biased/vector-separable generalized Good--Turing program listed above as open problems.

\appendix

\section{Algebra checklist}

This appendix collects the algebraic identities used in the proof.

\subsection{Derivative calculation}

For
\[
  \varphi_\theta(z)=z^\theta(z+\theta)e^{-z},
\]
one has
\begin{align*}
  \varphi_\theta'(z)
  &=\theta z^{\theta-1}(z+\theta)e^{-z}+z^\theta e^{-z}-z^\theta(z+\theta)e^{-z} \\
  &=z^{\theta-1}e^{-z}\{\theta(z+\theta)+z-z(z+\theta)\} \\
  &=z^{\theta-1}e^{-z}(\theta^2+z-z^2).
\end{align*}
Differentiating once more yields
\[
  \varphi_\theta''(z)=z^{\theta-2}e^{-z}
  \left[z^3-(\theta+2)z^2-\theta(\theta-1)z+
  \theta^2(\theta-1)\right].
\]

\subsection{\texorpdfstring{Curvature at the stationary point of $\varphi_\theta$}{Curvature at the stationary point of phi theta}}

At
\[
  \rho_\theta=\frac{1+\sqrt{1+4\theta^2}}{2},
  \qquad \rho_\theta^2=\rho_\theta+\theta^2,
\]
we have
\begin{align*}
  P_\theta(\rho_\theta)
  &=\rho_\theta^3-(\theta+2)\rho_\theta^2
    -\theta(\theta-1)\rho_\theta+\theta^2(\theta-1) \\
  &=\rho_\theta(1+\theta^2)+\theta^2
    -(\theta+2)(\rho_\theta+\theta^2)
    -\theta(\theta-1)\rho_\theta+\theta^2(\theta-1) \\
  &=-\rho_\theta-2\theta^2<0.
\end{align*}
This verifies that the two positive zeros of $P_\theta$ straddle $\rho_\theta$.

\begin{thebibliography}{99}

\bibitem{Good1953}
I. J. Good.
\newblock The population frequencies of species and the estimation of population parameters.
\newblock \emph{Biometrika}, 40(3--4):237--264, 1953.

\bibitem{McAllesterSchapire2000}
D. A. McAllester and R. E. Schapire.
\newblock On the convergence rate of Good--Turing estimators.
\newblock In \emph{Proceedings of the Thirteenth Annual Conference on Computational Learning Theory}, pages 1--6, 2000.

\bibitem{BerendKontorovich2012}
D. Berend and A. Kontorovich.
\newblock The missing mass problem.
\newblock \emph{Statistics and Probability Letters}, 82(6):1102--1110, 2012.

\bibitem{Painsky2023}
A. Painsky.
\newblock Generalized Good--Turing improves missing mass estimation.
\newblock \emph{Journal of the American Statistical Association}, 118(543):1890--1899, July 2023.
\newblock doi:10.1080/01621459.2021.2020658.

\bibitem{Painsky2022}
A. Painsky.
\newblock Convergence guarantees for the Good--Turing estimator.
\newblock \emph{Journal of Machine Learning Research}, 23(279):1--37, 2022.

\bibitem{FavaroNaulet2023}
S. Favaro and Z. Naulet.
\newblock Optimal estimation of high-order missing masses, and the rare-type match problem.
\newblock arXiv:2306.14998, 2023.

\end{thebibliography}

\end{document}
